% ════════════════════════════════════════════════════════════════
\subsection{Matroidy}
% ════════════════════════════════════════════════════════════════

\begin{definition} \label{def:matroid}
	\uemph{Matroid} to para $M = (S, \mathcal{I})$\footnote{Oznaczenia za CLRS: $S$ (\emph{set}) -- zbiór bazowy, $\mathcal{I}$ (\emph{independent}) -- rodzina zbiorów niezależnych. Uwaga: ,,niezależność'' jest tu cechą \emph{pojedynczego} zbioru ($A \in \mathcal{I}$ znaczy, że $A$ jest sam w~sobie niezależny), a~nie relacją między zbiorami -- dziedziczność (wł.\ \ref{matroid:dziedziczność}) mówi tylko, że podzbiór zbioru niezależnego też jest niezależny.} taka, że:
	\begin{enumerate}
		\item \label{matroid:skończoność} $S$ -- niepusty zbiór skończony,
		\item \label{matroid:niepustość} $\mathcal{I} \subseteq 2^S$ i $\mathcal{I} \neq \emptyset$ -- czyli $\mathcal{I}$ to niepusta rodzina podzbiorów $S$,
		\item \label{matroid:dziedziczność} $\forall A, B \subseteq S \quad A \in \mathcal{I} \wedge B \subseteq A \implies B \in \mathcal{I}$,
		\item \label{matroid:wymiana} własność wymiany:
		      \begin{flalign*}
			       & \forall A, B \subseteq S \quad A, B \in \mathcal{I} \wedge |B| > |A| \implies \text{istnieje } x \in B \setminus A \text{ taki, że } A \cup \{x\} \in \mathcal{I}. &  &  &  &
		      \end{flalign*}
	\end{enumerate}
	$\mathcal{I}$ -- rodzina zbiorów niezależnych matroidu $M$, elementy $\mathcal{I}$ -- zbiory niezależne matroidu $M$.

	Zbiór niezależny $A \in \mathcal{I}$ nazywamy \uemph{maksymalnym} (lub \uemph{bazą} matroidu $M$), jeśli nie da się go powiększyć, pozostając niezależnym, tzn.\ dla każdego $x \in S \setminus A$ zachodzi $A \cup \{x\} \notin \mathcal{I}$.
\end{definition}

\begin{dygresja}
	Matroid to abstrakcja i \emph{uogólnienie} pojęcia liniowej niezależności z~przestrzeni wektorowych. Aksjomaty są dokładnie tymi własnościami niezależności, które przetrwają, gdy zapomnimy o~samych wektorach i~zostawimy jedynie rodzinę ,,dozwolonych'' (niezależnych) podzbiorów: dziedziczność (wł.\ \ref{matroid:dziedziczność}) odpowiada temu, że podzbiór wektorów liniowo niezależnych nadal jest niezależny, a~własność wymiany (wł.\ \ref{matroid:wymiana}) to nic innego jak \uemph{lemat Steinitza o~wymianie}: jeśli $\{u_1, \dots, u_m\}$ jest liniowo niezależny i~leży w~przestrzeni rozpiętej przez $\{w_1, \dots, w_n\}$, to $m \leq n$ oraz można wymienić $m$ spośród wektorów $w_i$ na $u_1, \dots, u_m$, nie zmieniając rozpiętej przestrzeni --- mniejszy zbiór niezależny zawsze da się ,,dopełnić'' wektorami z~większego rozpinającego. To właśnie ta zdolność dokładania elementów uogólnia się do własności wymiany matroidu. Stąd słownictwo teorii matroidów --- ,,niezależny'', ,,baza'', ,,rząd'' --- pożyczone wprost z~algebry liniowej (oraz z~teorii grafów, por.\ matroid grafowy).

	Przykład \ref{ex:matroid:wektorowy} (matroid wektorowy) jest pierwowzorem; przykład \ref{ex:matroid:grafowy} pokazuje, że ta sama struktura opisuje też acykliczność w~grafach --- te dwa źródła, algebra liniowa i~teoria grafów, dały początek całej teorii.
\end{dygresja}

\begin{example} \label{ex:matroid-podstawowe} \leavevmode
	\begin{enumerate}[ref=\thedefinition.\arabic*]
		\item \label{ex:matroid:wektorowy} $A$ -- macierz, $S$ -- zbiór wierszy $A$, $\mathcal{I} = \{R \subseteq S : R \text{ tworzy zbiór liniowo niezależny}\}$.\\
		      Wtedy $(S, \mathcal{I})$ to matroid.
		\item \label{ex:matroid:grafowy} $G = (V, E)$ -- graf, $\mathcal{I} = \{F \subseteq E : G[F] \text{ jest acykliczny}\}$.\\
		      Wtedy $(E, \mathcal{I})$ to matroid zwany matroidem grafowym.
	\end{enumerate}
	Z definicji $\emptyset \in \mathcal{I}$.
\end{example}

\begin{dygresja}
	Wspólna intuicja obu przykładów: $S$ to zbiór ,,kawałków'', a $\mathcal{I}$ wyróżnia te podzbiory, w~których \emph{nie ma nadmiarowości}. Dla macierzy nadmiarowość to liniowa zależność wierszy, dla grafu --- cykl (krawędź, którą można usunąć bez rozspójniania tego, co już połączone). Dziedziczność (wł.\ \ref{matroid:dziedziczność}) mówi, że usunięcie kawałka z~niezależnego zbioru nie tworzy nadmiarowości; własność wymiany (wł.\ \ref{matroid:wymiana}) --- że mniejszy zbiór niezależny zawsze da się powiększyć o~element z~większego.

	Dla matroidu grafowego: zbiór krawędzi jest niezależny dokładnie wtedy, gdy tworzy las (jest acykliczny). Po lewej --- niezależny zbiór (las, krawędzie wyróżnione); po prawej --- ten sam zbiór z~dołożoną krawędzią $e_3$ zamykającą cykl $a\,b\,c$, więc już zależny:
	\begin{azfigure}
		\centering
		\begin{tikzpicture}[every node/.style={circle, draw, minimum size=0.6cm, inner sep=1pt}, baseline]
			% lewa: las (niezależny)
			\begin{scope}
				\node (a) at (0, 1.4) {$a$};
				\node (b) at (1.2, 1.4) {$b$};
				\node (c) at (0.6, 0.3) {$c$};
				\node (d) at (1.8, 0.3) {$d$};
				\draw[accent, very thick] (a) -- (b) node[draw=none, midway, above] {$e_1$};
				\draw[accent, very thick] (b) -- (c) node[draw=none, midway, right=1pt] {$e_2$};
				\draw[accent, very thick] (c) -- (d) node[draw=none, midway, below] {$e_4$};
				\node[draw=none] at (0.9, -0.6) {$F = \{e_1, e_2, e_4\} \in \mathcal{I}$};
			\end{scope}
			% prawa: cykl (zależny)
			\begin{scope}[xshift=5cm]
				\node (a) at (0, 1.4) {$a$};
				\node (b) at (1.2, 1.4) {$b$};
				\node (c) at (0.6, 0.3) {$c$};
				\node (d) at (1.8, 0.3) {$d$};
				\draw[accent, very thick] (a) -- (b) node[draw=none, midway, above] {$e_1$};
				\draw[accent, very thick] (b) -- (c) node[draw=none, midway, right=1pt] {$e_2$};
				\draw[accent, very thick] (c) -- (d) node[draw=none, midway, below] {$e_4$};
				\draw[accentB, very thick] (a) -- (c) node[draw=none, midway, left=1pt] {$e_3$};
				\node[draw=none] at (0.9, -0.6) {$F \cup \{e_3\} \notin \mathcal{I}$};
			\end{scope}
		\end{tikzpicture}
	\end{azfigure}

	Ten sam obraz dla matroidu wektorowego (przykład \ref{ex:matroid:wektorowy}) w~$\R^2$: $S = \{v_1, v_2, v_3\}$, niezależność = liniowa niezależność. Dwa niewspółliniowe wektory są niezależne, ale dowolne trzy wektory w~$\R^2$ są już zależne --- każdy zbiór $\geq 3$ elementów jest zależny, więc bazami są wszystkie pary niewspółliniowe (rząd $2$):
	\begin{azfigure}
		\centering
		\begin{tikzpicture}[>=stealth, baseline]
			% lewa: para niezależna
			\begin{scope}
				\draw[->, accent, very thick] (0,0) -- (1.6, 0.3) node[above right] {$v_1$};
				\draw[->, accent, very thick] (0,0) -- (0.5, 1.5) node[above] {$v_2$};
				\node at (1.0, -0.8) {$\{v_1, v_2\} \in \mathcal{I}$};
			\end{scope}
			% prawa: trzy wektory zależne
			\begin{scope}[xshift=5cm]
				\draw[->, accent, very thick] (0,0) -- (1.6, 0.3) node[above right] {$v_1$};
				\draw[->, accent, very thick] (0,0) -- (0.5, 1.5) node[above] {$v_2$};
				\draw[->, accentB, very thick] (0,0) -- (1.5, 1.2) node[right] {$v_3$};
				\node at (1.0, -0.8) {$\{v_1, v_2, v_3\} \notin \mathcal{I}$};
			\end{scope}
		\end{tikzpicture}
	\end{azfigure}
\end{dygresja}

\begin{lemma} \label{lemma:matroid:liczność}
	Niech $M = (S, \mathcal{I})$. Wszystkie maksymalne zbiory niezależne $M$ mają tę samą liczność.
\end{lemma}
\begin{proof}
	Niech $A, B \in \mathcal{I}$ będą maksymalnymi zbiorami niezależnymi $M$. Przypuśćmy, że $|B| > |A|$.\\
	Z własności \ref{matroid:wymiana}. istnieje $x \in B \setminus A$ taki, że $A \cup \{x\} \in \mathcal{I}$ -- sprzeczność z maksymalnością $A$.
\end{proof}

\begin{problem}
$M = (S, \mathcal{I})$ -- matroid, $w: S \to \Rp$ -- funkcja wagi.\\
Dla $A \subseteq S$ niech $w(A) = \sum_{x \in A} w(x)$ -- waga $A$.\\
Znaleźć zbiór niezależny o maksymalnej wadze w $M$.\\
$S = \{x_1, x_2, \ldots, x_n\}$.
\end{problem}

\begin{alg}{Algorytm Greedy}{alg:greedy}
	\FunctionNN{Greedy}{$M, w$}
	\State $n \Assign |S|$
	\State $A \Assign \emptyset$ \Comment{$A$ -- rozwiązanie}
	\State posortuj $S$ w porządku nierosnącym względem wagi, tak żeby $w(x_1) \geq w(x_2) \geq \ldots \geq w(x_n)$
	\For{$i \Assign 1$ \To $n$}
	\If{$A \cup \{x_i\} \in \mathcal{I}$} \label{line:greedy-niezaleznosc}
	\State $A \Assign A \cup \{x_i\}$
	\EndIf
	\EndFor
	\State \Return $A$
	\EndFunction
\end{alg}

\begin{dygresja}
	Idea jest możliwie najprostsza: skoro chcemy dużej wagi, bierzmy elementy od najcięższego do najlżejszego i~dokładajmy każdy z~nich do rozwiązania, o~ile tylko nie psuje to niezależności (warunek z~linii \ref{line:greedy-niezaleznosc}). Raz dodanego elementu nigdy nie wycofujemy --- stąd nazwa \emph{zachłanny}: na każdym kroku podejmujemy lokalnie najlepszą decyzję i~się jej trzymamy, nie martwiąc się o~przyszłość.

	Dla dowolnej rodziny $\mathcal{I}$ taka krótkowzroczność zwykle zawodzi, bo wczesny ciężki wybór może zablokować dwa lepsze późniejsze. To, co ratuje algorytm tutaj, to struktura matroidu. Własność wymiany (wł.\ \ref{matroid:wymiana}) gwarantuje, że zachłanny wybór nigdy nas nie ,,zamknie'' w~gorszym rozwiązaniu: jeśli istnieje cięższy zbiór niezależny, to zawsze da się go wykorzystać do ulepszenia bieżącego --- dokładnie ten argument napędza dowód twierdzenia \ref{thm:rado-edmonds}. Co więcej, działa to też w~drugą stronę: zachłanność daje optimum dla \uemph{każdej} funkcji wagi \emph{wtedy i~tylko wtedy}, gdy rodzina $\mathcal{I}$ jest matroidem --- to właśnie czyni matroidy ,,naturalnym'' środowiskiem dla algorytmów zachłannych.
\end{dygresja}

\begin{theorem}[Rado, Edmonds] \label{thm:rado-edmonds}
	Algorytm Greedy znajduje zbiór niezależny o maksymalnej wadze.
\end{theorem}
\begin{proof}
	Niech $A = \{x_{i_1}, x_{i_2}, \ldots, x_{i_k}\}$ będzie zbiorem zwróconym przez algorytm, przy czym $w(x_{i_1}) \geq w(x_{i_2}) \geq \ldots \geq w(x_{i_k})$. Dowód przebiega w~dwóch krokach: najpierw pokazujemy, że $A$ jest maksymalny, a~następnie, że żaden zbiór niezależny nie ma większej wagi.

	\emph{Krok 1: $A$ jest maksymalnym zbiorem niezależnym.}
	Przypuśćmy przeciwnie, że istnieje $x_j \in S \setminus A$ taki, że $A \cup \{x_j\} \in \mathcal{I}$. Wtedy $j < i_k$, bo w~przeciwnym wypadku algorytm dodałby $x_j$ do $A$ jako $(k+1)$-szy element. Niech $l$ będzie najmniejszą liczbą taką, że $j < i_l$; wówczas
	\[
		w(x_{i_1}) \geq \ldots \geq w(x_{i_{l-1}}) \geq w(x_j) \geq w(x_{i_l}).
	\]
	Ponieważ $A \cup \{x_j\} \in \mathcal{I}$, z~własności \ref{matroid:dziedziczność} zbiór $\{x_{i_1}, \ldots, x_{i_{l-1}}, x_j\} \subseteq A \cup \{x_j\}$ jest niezależny. Ale wtedy algorytm dodałby $x_j$ zamiast $x_{i_l}$ jako $l$-ty element --- sprzeczność. Zatem $A$ jest maksymalnym zbiorem niezależnym.

	\emph{Krok 2: $w(B) \leq w(A)$ dla każdego zbioru niezależnego $B$.}
	Rozważmy zbiór niezależny $B = \{y_1, y_2, \ldots, y_m\} \in \mathcal{I}$ z~$w(y_1) \geq w(y_2) \geq \ldots \geq w(y_m)$. Z~maksymalności $A$ (Lemat \ref{lemma:matroid:liczność}) mamy $m \leq k$. Pokażemy, że
	\[
		w(y_j) \leq w(x_{i_j}) \qquad \text{dla każdego } j \leq m.
	\]
	Przypuśćmy przeciwnie, że $w(y_j) > w(x_{i_j})$ dla pewnego $j$, i~rozważmy zbiory niezależne $A_{j-1} = \{x_{i_1}, \ldots, x_{i_{j-1}}\} \subseteq A$ oraz $B_j = \{y_1, \ldots, y_j\} \subseteq B$. Ponieważ $|B_j| = j > j-1 = |A_{j-1}|$, z~własności \ref{matroid:wymiana} istnieje $y_s \in B_j \setminus A_{j-1}$ taki, że $A_{j-1} \cup \{y_s\} \in \mathcal{I}$. Ponieważ $y_s \in B_j$, a~wagi $B$ są nierosnące, mamy $w(y_s) \geq w(y_j)$, skąd $w(y_s) \geq w(y_j) > w(x_{i_j})$; istnieje więc $t \leq j$ takie, że
	\[
		w(x_{i_1}) \geq \ldots \geq w(x_{i_{t-1}}) \geq w(y_s) > w(x_{i_t}).
	\]
	Z~własności \ref{matroid:dziedziczność} zbiór $\{x_{i_1}, \ldots, x_{i_{t-1}}, y_s\} \subseteq A_{j-1} \cup \{y_s\}$ jest niezależny. Ale wtedy algorytm dodałby $y_s$ zamiast $x_{i_t}$ jako $t$-ty element --- sprzeczność.

	Sumując nierówności po wszystkich $j \leq m$, otrzymujemy
	\[
		w(B) = \sum_{j=1}^{m} w(y_j) \leq \sum_{j=1}^{m} w(x_{i_j}) \leq w(A),
	\]
	co kończy dowód.
\end{proof}

\paragraph{Złożoność czasowa algorytmu Greedy}
$n = |S|$
\begin{itemize}[nosep]
	\item linia 1: $\BO(1)$
	\item linia 2: $\BO(1)$
	\item linia 3: $\BO(n \log n)$
	\item liczba iteracji pętli z linii 4 wynosi $n$
	\item linia 5: $f(n)$ -- czas sprawdzenia warunku z linii 5
	\item linia 6: $\BO(1)$
	\item linia 7: $\BO(1)$
\end{itemize}
Pętla wykonuje $n$ iteracji, a~każda kosztuje $f(n) + \BO(1)$, co daje koszt pętli $\BO\bigl(n \cdot f(n)\bigr)$. Łącznie z~sortowaniem otrzymujemy
\[
	\BO\bigl(n \log n + n \cdot f(n)\bigr),
\]
gdzie $f(n)$ to czas sprawdzenia niezależności $A \cup \{x_i\} \in \mathcal{I}$ (linia \ref{line:greedy-niezaleznosc}). Złożoność zależy więc od reprezentacji matroidu: dla matroidów, w~których test niezależności jest wielomianowy, cały algorytm działa w~czasie wielomianowym.

% ════════════════════════════════════════════════════════════════
\subsection{Problem szeregowania zadań}
% ════════════════════════════════════════════════════════════════

\begin{itemize}[nosep]
	\item $S = \{z_1, z_2, \ldots, z_n\}$ -- zbiór zadań o jednostkowym czasie wykonania
	\item $d_1, d_2, \ldots, d_n$ -- deadliny, $d_i \in \{1, \ldots, n\}$ dla $i \in [n]$
	\item $w_1, w_2, \ldots, w_n$ -- kary, $w_i \geq 0$ dla $i \in [n]$
\end{itemize}
Płacimy karę $w_i$ za zadanie $z_i$, jeżeli nie jest skończone w momencie $d_i$.
\begin{problem}[szeregowania zadań z karami]
Znaleźć kolejność wykonania zadań minimalizującą sumę kar.
\end{problem}

Załóżmy, że mamy pewien harmonogram. Zadanie $z_i$ nazywamy \uemph{terminowym} w tym harmonogramie, jeśli jest zakończone do chwili $d_i$, w przeciwnym razie nazywamy je \uemph{spóźnionym}.\\
Nie płacimy kar za zadania terminowe, płacimy kary za zadania spóźnione.

\begin{lemma}
	Każdy harmonogram można przetransformować na harmonogram z taką samą sumą kar, w którym wszystkie zadania terminowe poprzedzają wszystkie zadania spóźnione oraz zadania terminowe są posortowane w kolejności niemalejącej względem deadlinów.
\end{lemma}
\begin{proof}~
	\begin{enumerate}
		\item Niech $z_i$ będzie zadaniem spóźnionym ukończonym w chwili $l_i > d_i$ i niech $z_j$ będzie zadaniem terminowym ukończonym w chwili $l_j \leq d_j$ i niech $z_i$ będzie przed $z_j$ w harmonogramie, tzn. $l_i < l_j$.\\
		      Zamieniamy miejscami $z_i$ i $z_j$. Po zamianie:
		      \begin{itemize}[nosep]
			      \item $z_i$ jest ukończone w chwili $l_j > l_i > d_i$, więc nadal jest spóźnione.
			      \item $z_j$ jest ukończone w chwili $l_i < l_j \leq d_j$, więc nadal jest terminowe.
		      \end{itemize}
		      Po pewnej liczbie takich zamian otrzymujemy harmonogram o tej samej sumie kar, w którym wszystkie zadania terminowe poprzedzają wszystkie zadania spóźnione.

		      \begin{azfigure}
			      \centering
			      \begin{tikzpicture}[
					      slot/.style={rectangle, draw, minimum width=0.95cm, minimum height=0.7cm, inner sep=1pt, font=\small},
					      late/.style={slot, draw=accentB, fill=accentB!18},
					      ontime/.style={slot, draw=accent, fill=accent!12},
					      empty/.style={slot, draw=gray!45, text=gray!70},
					      lab/.style={anchor=east, font=\small\itshape},
				      ]
				      % przed zamianą
				      \node[lab] at (0.4,1.3) {przed:};
				      \node[empty]  at (1,1.3) {};
				      \node[late]   (ti) at (2,1.3) {$z_i$};
				      \node[empty]  at (3,1.3) {};
				      \node[empty]  at (4,1.3) {};
				      \node[ontime] (tj) at (5,1.3) {$z_j$};
				      \node[empty]  at (6,1.3) {};
				      % po zamianie
				      \node[lab] at (0.4,0) {po:};
				      \node[empty]  at (1,0) {};
				      \node[ontime] (bj) at (2,0) {$z_j$};
				      \node[empty]  at (3,0) {};
				      \node[empty]  at (4,0) {};
				      \node[late]   (bi) at (5,0) {$z_i$};
				      \node[empty]  at (6,0) {};
				      % strzałki zamiany
				      \draw[->, accentB, thick] (ti) to[bend right=18] (bi);
				      \draw[->, accent,  thick] (tj) to[bend left=18]  (bj);
				      % oś czasu
				      \draw[->, gray] (0.6,-0.65) -- (6.5,-0.65) node[right, font=\small, text=gray] {czas};
			      \end{tikzpicture}
			      \caption*{Zadanie spóźnione $z_i$ (\textcolor{accentB}{pomarańczowe}) poprzedza zadanie terminowe $z_j$ (\textcolor{accent}{niebieskie}). Po zamianie $z_i$ kończy się jeszcze później (nadal spóźnione), a~$z_j$ wcześniej (nadal terminowe) --- suma kar bez zmian. Krok~2 to ten sam ruch między dwoma zadaniami terminowymi w~złej kolejności deadlinów.}
		      \end{azfigure}

		\item Niech $z_i, z_j$ będą zadaniami terminowymi takimi, że $z_i$ jest ukończone w chwili $l_i \leq d_i$, $z_j$ jest ukończone w chwili $l_j \leq d_j$, $l_i < l_j$, ale $d_i > d_j$.\\
		      Zamieniamy miejscami $z_i$ i $z_j$. Po zamianie:
		      \begin{itemize}[nosep]
			      \item $z_i$ jest ukończone w chwili $l_j \leq d_j < d_i$, więc $z_i$ jest nadal terminowe.
			      \item $z_j$ jest ukończone w chwili $l_i < l_j \leq d_j$, więc $z_j$ jest nadal terminowe.
		      \end{itemize}
		      Po pewnej liczbie takich zamian otrzymujemy harmonogram spełniający warunki lematu.
	\end{enumerate}
\end{proof}
